\section{政策节点下的统计检验与稳健性分析}

在描述统计和风险指数之外，本文使用固定效应回归考察政策节点后美国来源进口额的相对变化。基准模型设定为：
\[
\ln(Import_{i,p,t}+1)=\beta_1US_i\times Post2018_t+\beta_2US_i\times Post2022_t+\beta_3US_i\times Post2023_t+\alpha_i+\gamma_p+\delta_t+\varepsilon_{i,p,t}.
\]
其中 $\alpha_i$、$\gamma_p$ 和 $\delta_t$ 分别表示出口国固定效应、产品固定效应和年份固定效应。被解释变量首先采用进口额对数形式，并补充以进口份额为被解释变量的估计。

\begin{table}[H]
\centering
\caption{政策节点回归核心结果：进口额（OLS）与进口份额（PPML）}
\label{tab:regression}
\begin{tabular}{llrrc}
\toprule
 & 变量 & 系数 & 标准误 & p值 \\
\midrule
\multicolumn{5}{l}{\textit{列1：OLS，被解释变量 $\ln(Import+1)$}} \\
 & US$\times$Post2018 & 0.177 & 0.058 & 0.002 \\
 & US$\times$Post2022 & 0.026 & 0.048 & 0.589 \\
 & US$\times$Post2023 & -0.131 & 0.052 & 0.012 \\
\midrule
\multicolumn{5}{l}{\textit{列2：PPML，被解释变量 进口份额}} \\
 & US$\times$Post2018 & -0.053 & 0.074 & 0.478 \\
 & US$\times$Post2022 & -0.293 & 0.031 & $<$0.001 \\
 & US$\times$Post2023 & -0.320 & 0.062 & $<$0.001 \\
\bottomrule
\end{tabular}

\vspace{0.5em}
\begin{minipage}{0.92\linewidth}
\footnotesize
注：两列模型均包含出口国、产品和年份三重固定效应。列1使用按出口国聚类的稳健标准误，列2使用 PPML 估计，不报告传统 $R^2$。样本量 10\,608。
\end{minipage}
\end{table}


表 \ref{tab:regression} 显示，绝对进口额模型中 US$\times$Post2018、US$\times$Post2022 和 US$\times$Post2023 均不显著。份额型模型中三个交互项方向为负，但同样不显著。这一结果与描述统计并不矛盾：描述统计显示美国份额在部分产品中下降，但固定效应模型提醒我们，不能把这种结构变化简单解释为美国出口管制的严格因果效应。

稳健性检验包括改用反双曲正弦变换、剔除 2020--2021 年以及收窄到 2024 年前十来源国加美国的对照组。相关结果未改变核心判断：当前模型更适合支持“来源结构重组”和“美国份额承压”的经验事实，而不适合支持“美国对华出口额显著下降”的强因果结论。
